\section{Parameter constraints and attack costs}
\label{sec:sizing}
\begin{definition}[Sizing envelope]
An envelope $E$ fixes target ledger fee parameters, transaction and script sizes,
execution budgets, output-minimum rules, supported datum sizes, and assumed
inclusion and response delays. Let $c_x(E)$ bound the measured ledger fee of
a supported singleton exit $x$, and let $e_x(E)$ be a chosen bound on external
processor cost. Neither function is currently measured by this paper.
\end{definition}

\proposed If a processor requires nonnegative net income within $E$, an exit
with gross tip $t$ needs $t-c_x(E)-e_x(E)\ge0$. Consequently a sufficient
conditional floor is
\[
 t_{\min}\ \ge\ \max_{x\in X}\{c_x(E)+e_x(E)\},
\]
where $X$ contains only exits with an implemented, balanced allocation paying
that gross tip. The measured fee-only requirement is the special case $e_x=0$.
The current fold equation supports this arithmetic. Rejection permits a
processor to leave some tip to the requester, so a cost-cover construction
needs to choose its refund and fee allocation explicitly. Thus the formula is
not a completed all-exit implementation result. Retraction returns the tip
and is outside $X$: its checked model construction uses ordinary lovelace
funding $I\ge f$, supplied by the requester in this cost model. Increasing
its refundable tip alone does not pay the fee or a processor. Its ledger fee
still needs measurement in $E$.

\hand If measured fold fee were at most $170{,}000$ and external cost at most
$30{,}000$ lovelace, a $200{,}000$ floor would break even in that hypothetical
envelope. A later fee of $260{,}000$ would give $-90{,}000$ net with the same
external cost. No measurement or recommendation is attached to these numbers.

\proposed The design criterion is $0<\delta<W_p$. For desired booking
inclusion slack $d_b$, processing slack $d_p$ after booked time, and reclaim
response slack $d_r$, require $d_b\le\delta$, $d_p<W_p$ and $d_r<W_r$.
These are conditional feasibility requirements; no consensus inclusion bound
establishes $d_b,d_p,d_r$. The checked booking theorem enforces width at most
$\delta$; it does not establish this entire cross-parameter sizing criterion
for every decoded configuration.

\hand Given $\tau\in[L,U)$ and $s=U$, then $0<s-\tau\le\delta$.
Hence a compatible registration's expiry threshold is between
$\tau+W_p+W_r$ and $\tau+\delta+W_p+W_r$. Adding nonnegative windows to the
checked booking inequalities proves this bound. It is a threshold bound only:
a current mismatch is available earlier, and retraction needs an executable
fee-paying transaction as well as time.

\proposed To make an adversary pay at least a selected treasury charge $C$
for $n>0$ successful junk registrations, require $nR\ge C$, equivalently
$R\ge\lceil C/n\rceil$. This controls only donations of successful folds.
Selecting $C$ is a policy choice, not an optimal price or a universal flood
deterrent. A party assigning value $V(k)>R+t+b+a$ to a key can still choose
to buy it under these simple monetary assumptions.

\hand The following costs follow only under the expressly stated allocation
and ownership assumptions; none is a game-theoretic equilibrium claim.
Successful junk folding pays $R$ per key to the treasury, plus ledger fees;
if the funder and processor coincide, the tip is an internal transfer.
If that same party also controls the destination, its carrier remains an
asset, with an opportunity cost rather than an immediate wealth loss.
For a balanced exact-refund self-rejection, the tip is likewise internal;
the combined party loses the booking and rejection fees, plus external and
capital costs. There is no positive treasury floor for this attack.
For a losing race processed by someone else under the same exact-refund
assumption, the funder loses $t+b$ while receiving $R+a$ back. A larger tip
does not buy priority. A termination fold irrevocably donates its declared
carrier $a$; with no guaranteed delay bound, $rH\Delta$ remains unbounded.
For the funded retraction construction with all change returned to its
funder, that party loses $b+f$ plus external and capital costs; it recovers
the registration cost, carrier and tip. There is no treasury charge on that exit.

No bound relates the number of harmed requesters to the number of charged
attacker folds. A griefing ratio would divide total victim losses, including
delay costs, by total attacker expenditure under an explicit schedule. Neither
that schedule nor a finite delay bound is supplied; no finite universal ratio
is derived from a maximum batch size.

Policy identities and the registry identity are separation parameters, not
economic prices. They must identify the intended application and asset policies;
choosing them incorrectly cannot be repaired by a larger tip or window.
Carrier values must satisfy actual output minima when outputs are created.
They are not sized by a fixed registry constant. Changes outside $E$ require
new measurements and may motivate a new registry; no automatic migration or
rescue of old pending requests is implied.
